% ============================================================
% 04_theory.tex
% Section 4: Theoretical Reflection
% ============================================================

\section{Theoretical Reflection}\label{sec:theory}

Taken together, the theories discussed in this section explain the same challenge from different levels of analysis. The jagged frontier highlights the uneven capabilities of AI systems, knowledge codification explains how expertise is embedded into digital artifacts, and augmentation theory clarifies the continuing role of human judgment. Building on these foundations, the TOE framework helps explain why technically feasible systems still struggle to achieve organizational adoption. Together, they suggest that successful digital innovation depends not only on technical capability, but also on the organizational conditions required for adoption.

\subsection{AI and Work: Augmentation, Automation, and the Jagged Frontier}

\subsubsection{The Jagged Frontier in Practice}

Dell'Acqua et al.\parencite{DellAcqua2026} introduce the ``jagged technological frontier'' to describe how AI assistance improves performance on some tasks while degrading it on others, even within the same knowledge workflow and at similar perceived difficulty. Our build process provides a practical example of navigating this frontier.

\textbf{Inside the frontier} (AI helped):
\begin{itemize}
\item Boilerplate Python class definitions, docstrings, and type hints.
\item LaTeX equation formatting and notation tables.
\item Multi-language UI translation (200+ terms from Chinese to pharmaceutical-logistics English).
\item CSS animation and responsive layout code.
\end{itemize}

\textbf{Outside the frontier} (AI harmed or misled):
\begin{itemize}
\item Dependency management for Cloud deployment: the AI added \texttt{torchvision} to solve a symptom rather than the root cause.
\item Reward-function tuning: an initial $\alpha_{\text{setup}} = 5.0$ generated excessive small waves, a behavioral issue the AI failed to recognize.
\item Product positioning: the AI consistently expanded feature scope (NSGA-II, federated learning, 3D visualization) without evaluating whether users actually needed these additions.
\end{itemize}

The last point echoes Shipper's argument in ``After Automation''\parencite{shipper2026after}: as AI commoditizes codified expertise, the value of human judgment increases. While the AI could generate features rapidly, deciding which features created meaningful user value remained a fundamentally human task.

\subsubsection{Knowledge Codification and Its Risks}

Jurowetzki\parencite{Jurowetzki2026}, building on Cowan et al.\parencite{cowan2000codification}, decomposes codification into four capacities: encoding, evaluation, maintenance, and revision. AI dramatically lowers encoding costs by generating code and documentation, but it often increases maintenance and revision challenges.

Consider our \texttt{requirements.txt}. The AI encoded dependencies quickly. However, when the deployment environment changed---through updates to \texttt{transformers}, dependency conflicts, or Streamlit Cloud behavior---the responsibility for evaluating and revising the encoded knowledge fell entirely to the human developer.

This generates what Jurowetzki calls \textbf{orphaned codification}: artifacts whose assumptions are difficult to trace or challenge. Our project contains thousands of lines of AI-generated code whose underlying logic may become increasingly difficult to maintain as platforms evolve.

\subsubsection{Augmentation vs. Automation}

Raisch and Krakowski\parencite{raisch2021automation} describe the ``automation--augmentation paradox'': AI systems that automate routine tasks often increase demand for higher-level human judgment. Maverick-SORT was deliberately designed as an augmentation tool rather than an automation replacement.

Warehouse managers remain the decision-makers; the system recommends but does not execute. Even the What-If simulator is framed as a decision-support tool rather than an autonomous optimizer. This design limits immediate efficiency gains compared with full automation, but it substantially lowers adoption risk because users retain control and can override recommendations.

\subsection{Adoption and Diffusion: A TOE Framework Perspective}

The Technology--Organization--Environment (TOE) framework provides a useful lens for understanding whether Maverick-SORT can move beyond a prototype and achieve real operational adoption.

\paragraph{Technology.}
Maverick-SORT is designed as a decision-support layer rather than a replacement for existing WMS platforms. While WMS systems excel at transaction processing and execution, they provide limited support for simulation, optimization, and scenario analysis. By remaining standalone and browser-based, the system delivers predictive intelligence without requiring costly modifications to legacy infrastructure.

\paragraph{Organization.}
The key adopter is not the IT department but the Operations Director. Unlike IT managers, whose priority is stability and risk avoidance, operations managers are directly responsible for warehouse performance. The Essential Edition lowers adoption barriers by allowing them to evaluate potential value independently through a zero-deployment What-If environment.

\paragraph{Environment.}
Pharmaceutical logistics operates under strict GSP compliance requirements. Rather than treating regulation as an obstacle, Maverick-SORT uses explainable, knowledge-guided recommendations to align optimization with existing operational rules. Regulatory pressure therefore becomes a driver of adoption rather than a barrier.

\paragraph{The Most Limiting Condition.}
Among all TOE factors, the greatest constraint is organizational champion bandwidth. Operations managers often lack the time, attention, and political capital required to evaluate new technologies. Consequently, adoption depends less on algorithmic performance than on reducing the effort required to explore the solution. The Essential Edition, free trial, and five-minute simulation workflow are all designed to minimize this activation cost.

\subsection{Demo Day as an Adoption Test}

The class Demo Day provided a useful reality check regarding how the system was perceived compared with how it was built.

\begin{itemize}
\item \textbf{What resonated most (System Logic).} Reviewers were most engaged by the relationships between components---particularly how KGDRL recommendations connected to the Algorithm Arena and AI Copilot explanations.

```
\item \textbf{What was largely ignored (Visual Complexity).} Detailed dashboard visualizations attracted relatively little attention. Without sufficient operational context, many charts appeared as generic analytics rather than sources of business value.

\item \textbf{The core skepticism (Product Identity).} The strongest critiques questioned the product's fundamental value proposition: \textit{What is the core benefit? Why would someone adopt this system?}
```

\end{itemize}

\textbf{The Key Lesson.} This feedback revealed that my immediate challenge was not technical performance but narrative clarity. As a result, we shifted the project's positioning away from a feature-oriented demonstration and toward the concept of a ``Risk Reversal Engine'' designed to reduce adoption friction. The experience reinforced the broader conclusion of this report: organizational attention, rather than algorithm quality, is often the primary constraint on digital innovation adoption.
